% In dieser Umgebung wird die Titelseite separat vom restlichen Text gesetzt
\begin{titlepage}
  % andere Seitenränder als im Rest der Arbeit
  \newgeometry{lmargin=2cm,tmargin=7mm,rmargin=5mm,bmargin=1cm}
  % die "Hausfarbe" der HAW; diese und die folgenden Einstellungen sind lokal
  % und gelten nur innerhalb der Umgebung "titlepage"
  \color{haw}
  % Blocksatz für die Titelseite deaktivieren
  \raggedright
  % Logo rechtsbündig setzen
  \hfill\includegraphics[width=7cm]{HAW_Marke_RGB_300dpi}\\

  % vertikaler Abstand
  \vspace{5cm}

  % Wahl der "Hausschrift" Open Sans der HAW, die als Schrift auf Ihrem
  % Rechner installiert sein muss
  \setmainfont{Open Sans}
  % etwas kleiner als üblich
  \small
  % fett und in Majuskeln
  \textbf{BACHELORARBEIT}

  % vertikaler Abstand
  \vspace{8mm}

  % der Titel der Arbeit als "Seite in der Seite"; natürlich müssen Sie hier
  % Ihren Titel eintragen
  \begin{minipage}{0.8\linewidth}
    % Wahle der zweiten "Hausschrift" der HAW, die ebenfalls auf Ihrem Rechner
    % bereits vorhanden sein muss
    \setmainfont{Martel Heavy}
    % ziemlich große Schrift
    \LARGE
    % [1mm] steht jeweils für einen etwas größeren Durchschuss TODO Titel fixen
    CodeGreen:\\[1mm]
    Energieverbrauch von Websites\\[1mm]
    Tool zur Verbesserung\\[1mm]
    des Energieverbrauchs von Code\\
    % am Ende noch ein waagerechter Strich, das CD will es so...
    \,\rule{11mm}{1.2mm}
  \end{minipage}

  % vertikaler Abstand, überraschenderweise
  \vspace{1cm}

  % hier korrektes Datum und Ihren Namen eingeben TODO Datum fixen
  vorgelegt am 1. Juli 2024\\
  Leon Beitz

  % letzter vertikaler Abstand für heute
  \vspace{5cm}

  % noch eine "Seite in der Seite", etwas nach rechts geschoben
  \hspace*{37mm}
  \begin{minipage}{0.5\linewidth}
    % Namen und Titel der beiden Prüfer eintragen TODO Titel fixen
    \begin{tabular}{@{}ll}
      Erstprüferin: & Prof. Dr.-Ing. Lars Hamann\\[-.3mm]
      Zweitprüfer: & Prof. Dr. Bettina Buth \\
    \end{tabular}\\

    % noch ein horizontaler Strich
    \,\rule{9mm}{1mm}\\[1.5mm]

    \textbf{HOCHSCHULE FÜR ANGEWANDTE}\\
    \textbf{WISSENSCHAFTEN HAMBURG}\\
    Department Informatik\\
    Berliner Tor 7\\
    20099 Hamburg
  \end{minipage}
\end{titlepage}
% setzt die Geometrie wieder auf die Standardwerte zurück
\restoregeometry

% für die Seite mit dem Abstract keine Seitenzahl ausgeben
\thispagestyle{empty}

\section*{Zusammenfassung}

% Hier ersetzen Sie bitte die vorhandenen Texte durch Ihre eigenen
% Zusammenfassungen
Der erste Teil der Arbeit befasst sich mit der Ermittlung von Methoden wie 
der Energieverbrauch von Websites ermittelt und reduziert werden kann. 
Dabei wird sich hauptsächlich mit dem Quellcode, der auf Entwickler Seite geschrieben
wird, befasst. Es werden Methoden zur Reduzierung des Energieverbrauchs
evaluiert und wie diese auf Quellcode angewendet werden können.
Der zweite Teil der Arbeit befasst sich mit der Umsetzung der gefundenen Methoden,
um ein Tool zu entwickeln, welches Verbesserungsvorschläge für Quellcode vorschlägt.
Dieses Tool benutzt 'Abstract Syntax Trees', um den Quellcode zu parsen und Patterns
zu finden. Hier wird sich den Herausforderungen und Abwägungen beim Auffinden von Patterns
auf 'Abstract Syntax Trees' beschäftigt.
Der dritte und letzte Teil der Arbeit